% Part: lambda-calculus
% Chapter: church-rosser
% Section: parallel-beta-reduction

\documentclass[../../../include/open-logic-section]{subfiles}

\begin{document}

\olfileid{lam}{cr}{pb}

\olsection{समांतर $\beta$-न्यूनीकरण}

या विभागात \emph{समांतर $\beta$-न्यूनीकरण} ही संकल्पना
मांडू आणि तिला चर्च–रॉसर गुणधर्म आहे हे सिद्ध करू.

\begin{defn}[समांतर $\beta$-न्यूनीकरण, $\bredpar$] \ollabel{defn:bredpar}
  पदांचे समांतर न्यूनीकरण ($\bredpar$) पुढीलप्रमाणे
  विगमनाने परिभाषित केले जाते:
  \begin{enumerate}
    \item \ollabel{defn:bredpar1} $x \bredpar x$.
    \item \ollabel{defn:bredpar2} $N \bredpar N'$ असेल, तर $\lambd[x][N] \bredpar
      \lambd[x][N']$.
    \item \ollabel{defn:bredpar3} $P \bredpar P'$ आणि $Q \bredpar Q'$ असतील, तर $PQ \bredpar
      P'Q'$.
    \item \ollabel{defn:bredpar4} $N \bredpar N'$ आणि $Q \bredpar Q'$ असतील, तर
      $(\lambd[x][N])Q \bredpar \Subst{N'}{Q'}{x}$.
  \end{enumerate}
\end{defn}

समांतर $\beta$-न्यूनीकरणात एका पायरीत पदातील कितीही
रेडेक्सचे संकुचन करता येते. ते $\beta$-न्यूनीकरणापेक्षा
वेगळे आहे: मूळ पदात असलेल्या रेडेक्सचेच संकुचन त्या
पायरीत करता येते; समांतर $\beta$-न्यूनीकरणामुळे नव्याने निर्माण
झालेल्या रेडेक्सचे नाही. उदाहरणार्थ,
$(\lambd[f][fx])(\lambd[y][y])$ या पदाचे समांतर
$\beta$-न्यूनीकरण केल्यास तेच पद किंवा
$(\lambd[y][y])x$ मिळू शकते; एका पायरीत पुढे~$x$
मिळू शकत नाही. मात्र नेहमीच्या $\beta$-न्यूनीकरणाने
ते~$x$ मध्ये न्यूनीत होते. कारण संबंधित रेडेक्स समांतर
$\beta$-न्यूनीकरणाची पहिली पायरी झाल्यावरच निर्माण
होतो. दुसरी समांतर $\beta$-न्यूनीकरणाची पायरी
घेतल्यास~$x$ मिळते.

\begin{thm}\ollabel{thm:refl}
  $M \bredpar M$.
\end{thm}

\begin{proof}
  सरावासाठी.
\end{proof}

\begin{prob}
  \olref[lam][cr][pb]{thm:refl} सिद्ध करा.
\end{prob}

\begin{defn}[$\beta$-पूर्ण विकास]\ollabel{defn:bcd}
  $M$ चा \emph{$\beta$-पूर्ण विकास}~$\bcd{M}$ पुढीलप्रमाणे
  विगमनाने परिभाषित केला जातो:
  \begin{align}
    \bcd{x} &= x \ollabel{defn:bcd1} \\
    \bcd{(\lambd[x][N])} &= \lambd[x][\bcd{N}] \ollabel{defn:bcd2}\\
    \bcd{(PQ)} &= \bcd{P}\bcd{Q} && \text{जर $P$ हे $\lambd$-अमूर्तीकरण नसेल तर}
    \ollabel{defn:bcd3} \\
    \bcd{((\lambd[x][N])Q)} &= \Subst{\bcd{N}}{\bcd{Q}}{x} \ollabel{defn:bcd4}
  \end{align}
\end{defn}

पदाचा $\beta$-पूर्ण विकास म्हणजे नावाप्रमाणे
``संपूर्ण समांतर न्यूनीकरण''. समांतर $\beta$-न्यूनीकरणात
एखाद्या रेडेक्सचे संकुचन टाळता येते; पण पूर्ण विकासात
सर्व मूळ रेडेक्सचे संकुचन करावेच लागते. म्हणून
$(\lambd[f][fx])(\lambd[y][y])$ चा पूर्ण विकास तेच
पद नसून $(\lambd[y][y])x$ आहे.

\begin{editorial}
  या व्याख्येत अडचण अशी की ($\lambd$-)पदांवरील फलने
  पुनरावर्ती रीतीने कशी परिभाषित करायची हे अजून
  मांडलेले नाही. पुढे दुरुस्त करू.
\end{editorial}

\begin{lem}\ollabel{lem:comp}
  $M \bredpar M'$ आणि $R \bredpar R'$ असतील, तर $\Subst{M}{R}{y}
  \bredpar \Subst{M'}{R'}{y}$.
\end{lem}

\begin{proof}
%READERNOTE{SOL6-A125: दोन आदेशनांच्या क्रमाची अदलाबदल करताना आवश्यक ताजेपणाच्या अटी आणि आदेशनाचे संयोजनसूत्र स्पष्ट केले आहे. आधी दुरुस्त केलेला अमूर्तीकरणाच्या प्रकरणातील न्यूनीत आदेशित प्रतिनिधी जपला आहे.}
अमूर्तीकरणांची नावे आवश्यक त्या सर्व मुक्त चरांपासून
आणि आदेशित चरांपासून वेगळी निवडू. पुढील आदेशनसूत्र वापरू:
\[
\Subst{\Subst{N}{Q}{x}}{R}{y}
=\Subst{\Subst{N}{R}{y}}{\Subst{Q}{R}{y}}{x},
\qquad x\ne y,\quad x\notin\FV{R}.
\]
हे सूत्र पदाच्या रचनेवरील विगमनाने सिद्ध होते. चराचे नाव
$x$ असल्यास दोन्ही बाजू $\Subst{Q}{R}{y}$ होतात;
नाव $y$ असल्यास दोन्ही बाजू $R$ होतात, कारण $R$ मध्ये
$x$ मुक्त नाही. इतर चर तसेच राहतात. उपयोजनात दोन्ही
उपपदांवर विगमन गृहीतक लागू होते. अमूर्तीकरणात त्याचे
बद्ध नाव दोन्ही आदेशित पदांच्या मुक्त चरांपासून आणि
$x,y$ पासून वेगळे निवडल्याने तेच गृहीतक देहाला लागू
करता येते. याने वर्गपदांवरील सूत्र सिद्ध होते.
या सिद्धतेतील बाह्य अमूर्तीकरणाचे नाव $x$ हे $y$ पासून
आणि $R,R'$ यांच्या सर्व मुक्त चरांपासून वेगळे निवडू.
  $M \bredpar M'$ च्या !!{derivation} वर विगमन करू.
  \begin{enumerate}
    \item अखेरची पायरी \olref{defn:bredpar1} आहे: सरावासाठी.
    \item अखेरची पायरी \olref{defn:bredpar2} आहे: मग $M$ हे
      $\lambd[x][N]$ आणि $M'$ हे $\lambd[x][N']$ आहे,
      जेथे $N \bredpar N'$. आपल्याला
      $\Subst{(\lambd[x][N])}{R}{y} \bredpar
      \Subst{(\lambd[x][N'])}{R'}{y}$, म्हणजे
      $\lambd[x][\Subst{N}{R}{y}] \bredpar
      \lambd[x][\Subst{N'}{R'}{y}]$, हे सिद्ध करायचे आहे.
      ते \olref{defn:bredpar2} आणि विगमन गृहीतकावरून
      लगेच मिळते.
    \item अखेरची पायरी \olref{defn:bredpar3} आहे: सरावासाठी.
    \item अखेरची पायरी \olref{defn:bredpar4} आहे: $M$ हे
      $(\lambd[x][N])Q$ आणि $M'$ हे $\Subst{N'}{Q'}{x}$ आहे.
      आपल्याला $\Subst{((\lambd[x][N])Q)}{R}{y} \bredpar
      \Subst{\Subst{N'}{Q'}{x}}{R'}{y}$, म्हणजे
      $(\lambd[x][\Subst{N}{R}{y}])\Subst{Q}{R}{y} \bredpar
      \Subst{\Subst{N'}{R'}{y}}{\Subst{Q'}{R'}{y}}{x}$,
      हे सिद्ध करायचे आहे. वरील आदेशनसूत्राने या दोन
      उजव्या बाजू समान आहेत. मग \olref{defn:bredpar4}
      आणि दोन्ही पूर्वपक्षांना लागू केलेल्या विगमन गृहीतकावरून
      आवश्यक समांतर न्यूनीकरण मिळते.
  \end{enumerate}
\end{proof}

\begin{prob}
  \olref[lam][cr][pb]{lem:comp} ची सिद्धता पूर्ण करा.
\end{prob}

\begin{lem}\ollabel{lem:cont}
  $M \bredpar M'$ असेल, तर $M' \bredpar \bcd{M}$.
\end{lem}
\begin{proof}
  $M \bredpar M'$ च्या !!{derivation} वर विगमन करू.
  \begin{enumerate}
    \item अखेरचा नियम \olref{defn:bredpar1} आहे: सरावासाठी.
    \item अखेरचा नियम \olref{defn:bredpar2} आहे: $M$ हे
      $\lambd[x][N]$ आणि $M'$ हे $\lambd[x][N']$ आहे,
      जेथे $N \bredpar N'$. आपल्याला $\lambd[x][N'] \bredpar
      \bcd{(\lambd[x][N])}$, म्हणजे \olref{defn:bcd2} नुसार
      $\lambd[x][N'] \bredpar \lambd[x][\bcd{N}]$,
      दाखवायचे आहे. ते \olref{defn:bredpar2} आणि
      विगमन गृहीतकावरून मिळते.
    \item अखेरचा नियम \olref{defn:bredpar3} आहे: $M$ हे $PQ$
      आणि $M'$ हे $P'Q'$ आहे, जेथे $P$, $Q$, $P'$, $Q'$
      अशी पदे आहेत की $P \bredpar P'$ आणि $Q \bredpar Q'$.
      विगमन गृहीतकानुसार $P' \bredpar \bcd{P}$ आणि
      $Q' \bredpar \bcd{Q}$.
      \begin{enumerate}
        \item काही $x$ आणि $N$ साठी $P$ हे $\lambd[x][N]$
          असेल, तर काही $N'$ साठी $P'$ हे $\lambd[x][N']$
          असावे लागते आणि $N \bredpar N'$. विगमन गृहीतकानुसार
          $N' \bredpar \bcd{N}$ आणि $Q' \bredpar \bcd{Q}$.
          मग \olref{defn:bredpar4} नुसार
          $(\lambd[x][N'])Q' \bredpar
          \Subst{\bcd{N}}{\bcd{Q}}{x}$.
        \item $P$ हे $\lambd$-अमूर्तीकरण नसेल, तर
          \olref{defn:bredpar3} नुसार $P'Q' \bredpar
          \bcd{P}\bcd{Q}$; आणि \olref{defn:bcd3} नुसार
          उजवीकडील पद $\bcd{PQ}$ आहे.
      \end{enumerate}
    \item अखेरचा नियम \olref{defn:bredpar4} आहे: काही
      $x$, $N$, $Q$, $N'$ आणि~$Q'$ साठी $M$ हे
      $(\lambd[x][N])Q$ आणि $M'$ हे $\Subst{N'}{Q'}{x}$,
      जेथे $N \bredpar N'$ आणि $Q \bredpar Q'$.
      विगमन गृहीतकानुसार $N' \bredpar \bcd{N}$ आणि
      $Q' \bredpar \bcd{Q}$. \olref{lem:comp} नुसार
      $\Subst{N'}{Q'}{x} \bredpar \Subst{\bcd{N}}{\bcd{Q}}{x}$;
      उजवीकडील पद नेमके $\bcd{((\lambd[x][N])Q)}$ आहे.
  \end{enumerate}
\end{proof}

\begin{prob}
  \olref[lam][cr][pb]{lem:cont} ची सिद्धता पूर्ण करा.
\end{prob}

\begin{thm}\ollabel{thm:cr}
  $\bredpar$ ला चर्च–रॉसर गुणधर्म आहे.
\end{thm}

\begin{proof}
  \olref{lem:cont} वरून लगेच मिळते.
\end{proof}

\end{document}
